% !PRINTED 684
\weberpage{684}
%<<CONT
und wenn man in den Punkt $\mathfrak{P}$ geht:
\weq{10}{\left(\frac{\beta-\beta_0}{\alpha-\alpha_0}\frac{d\alpha}{d\beta}\right)_0 = \frac{r}{s},}
was ein endlicher von Null verschiedener Wert ist.

Daraus folgt:

5. \gesperrt{Die Ordnungszahl des Differentialquotienten $d\alpha : d\beta$ in irgend einem Punkt $\mathfrak{P}$ ist gleich der Differenz der Ordnungszahlen von $\alpha - \alpha_0$ und $\beta - \beta_0$.}

Ist $r > s$, oder $r < s$, so ist
\[
\left(\frac{\alpha-\alpha_0}{\beta-\beta_0}\right)_0 = 0 \text{ oder} = \infty,
\]
und folglich nach (10) auch
\[
\left(\frac{d\alpha}{d\beta}\right)_0 = 0 \text{ oder} = \infty.
\]

Ist $r = s$, so folgt aus (10)
\weq{11}{\left(\frac{d\alpha}{d\beta}\right)_0 = \left(\frac{\alpha-\alpha_0}{\beta-\beta_0}\right)_0.}

\gesperrt{Hier ist jede Spur der Variablen $z$ herausgefallen, und diese Formel gilt für jeden Punkt $\mathfrak{P}$ ohne Ausnahme.}

Danach lassen sich die Nullpunkte und Unendlichkeitspunkte der Differentialquotienten genau feststellen und damit diese Funktionen als \gesperrt{Polygonquotienten} darstellen.

Ist $a$ die Ordnungszahl von $\alpha - \alpha_0$ in einem Punkt $\mathfrak{P}$ nach \S\,188, so enthält das Verzweigungspolygon $\mathfrak{Z}_\alpha$ den Punkt $\mathfrak{P}$ $a - 1$mal oder $-a-1$mal, je nachdem $a$ positiv oder negativ ist. Ist $a$ negativ, so enthält das Nennerpolygon $\mathfrak{A}$ von $\alpha$ den Punkt $\mathfrak{P}$ $(-a)$mal, und folglich ist in beiden Fällen der Punkt $\mathfrak{P}$ in $\mathfrak{Z}_\alpha : \mathfrak{A}^2$ $(a-1)$mal enthalten. Hat $\mathfrak{Z}_\beta, \mathfrak{B}, b$ die entsprechende Bedeutung für $\beta$, so enthält also der Quotient:
\[
\frac{\mathfrak{Z}_\alpha \mathfrak{B}^2}{\mathfrak{Z}_\beta \mathfrak{A}^2}
\]
$(a-b)$mal den Punkt $\mathfrak{P}$.

Die Ordnungszahl von $d\alpha : d\beta$ in $\mathfrak{P}$ ist aber nach 5. ebenso groß, und daraus ergibt sich die Darstellung:
\weq{12}{\frac{d\alpha}{d\beta} = \frac{\mathfrak{Z}_\alpha \mathfrak{B}^2}{\mathfrak{Z}_\beta \mathfrak{A}^2},}
worin, um es zu wiederholen, $\mathfrak{Z}_\alpha, \mathfrak{Z}_\beta$ die Verzweigungspolygone, $\mathfrak{A}, \mathfrak{B}$ die Nenner von $\alpha, \beta$ sind.
